Mit den Dichten $\rho_\text{Au} = \rAuO$ und $\rho_\text{Cu} = \rCuO$ (FoTaBe S.~187) und der Annahme, dass sich die Volumina der beiden Metalle addieren, ist
\begin{align}
 \rho &= \frac{m_\text{Au}+m_\text{Cu}}{V_\text{Au}+V_\text{Cu}} = \rLeF \\
   &= \frac{\mAuO+\mCuO}{\frac{\mAuO}{\rAuO}+\frac{\mCuO}{\rCuO}} \\
   &= \rLeg \approx \result{\rLegP{0}{3}}
\end{align}
Der gemessene Wert \rMeO\ ist grösser. Die Legierung hat also ein kleineres Volumen als Gold und Kupfer getrennt: Im gemeinsamen Kristallgitter sind die Atome dichter gepackt, die Volumina addieren sich nicht.
